Small lesions in brain magnetic resonance imaging (MRI) are difficult to segment because they occupy only a small fraction of the image volume and are dominated by background and larger lesions during voxel-wise optimization. A model can therefore reach a high Dice similarity coefficient (DSC) while missing many small lesions. We propose CATMIL, a training objective that augments the standard nnU-Net Dice and cross-entropy loss with two auxiliary terms and leaves the architecture unchanged. The Component-Adaptive Tversky (CAT) term weights lesion voxels by the inverse size of their connected component, so that each lesion receives a total weight and gradient nearly independent of its volume. The lesion-level Multiple Instance Learning (MIL) term treats each lesion as a bag of voxels and penalizes any lesion in which no voxel is detected. We evaluate CATMIL for multiple sclerosis lesion segmentation on MSLesSeg and on a second dataset, 3D-MR-MS, trained with the same loss weights. On MSLesSeg, CATMIL achieves the highest small-lesion recall (0.873 vs.\ 0.796 for Dice+CE; difference +0.077, 95\% CI +0.030 to +0.157, higher in all six test patients) and about 48\% fewer missed lesions, with DSC and HD95 comparable to the compared losses. The gain is a net gain in detected lesions, concentrated in lesions of at most 50~mm$^3$, and it extends to lesions of at least 3~mm in longest diameter, the size used in clinical reading (recall 0.944 vs.\ 0.870). Analysis of the predicted probabilities shows that standard losses produce no response to most of the small lesions they miss, so no decision threshold can recover them. The cost of the increased sensitivity is a larger number of small false-positive components, which lowers lesion-wise F1; a simple component-size filter removes most of them while preserving the sensitivity gain, and at the lesion-wise precision of Dice+CE, CATMIL detects more small lesions and has a higher lesion-wise F1. An ablation shows that the lesion-level MIL term accounts for the detection gain. On 3D-MR-MS, CATMIL again improves small-lesion recall and reduces missed lesions, with a larger false-positive cost and slightly lower DSC. All code is available at \url{https://github.com/luumsk/SmallLesionMRI}.
